\subsection{研究目的與取徑}\label{subsec:design}

本文的目的在把上述有方向的偏誤由一個須以模擬觀察的現象，轉為一個可以直接
計算的量，並據以設計校正。所循的路徑如下。標準 LC 的估計是對
$\log\bigl(\max(D_{x,t},c_0)/E_{x,t}\bigr)$ 作最小平方，
在 $M$-估計的架構中即取平方損失的估計量；而 $M$-估計量要在真值處一致，
其估計方程在真值處的期望須為零，該期望在反應變數先經非線性轉換之後不再
自動為零。本文因此計算標準 LC 的估計方程在真值處的期望，寫出其一般形式，
再檢視該形式在 Poisson、常態與過度離散三種分布下的具體樣貌。

此一路徑的兩項後果決定了全文的結構。偏誤的量既然可以寫出，
它就可以在配適之前算出，故本文的第一條主軸是\textbf{診斷}：
標準 LC 在某一組資料上會偏多少，由曝露數與一組粗略死亡率即可定出，
不需要模擬、真值或參考母體。偏誤的量既然可以算出，它也就可以扣除，
故第二條主軸是\textbf{校正}：扣除的位置決定了年齡水準與年齡敏感度兩項改善
得以疊加的條件。兩條主軸的適用範圍不同，前者只依賴分布假設，
後者另依賴秩一雙線性結構，此一差別於第\ref{sec:disc}節一併交代。

